% Fixed135 standalone substantive section.
% A 3-power tame radical tower; uniform odd-index Fox coefficient
% matrix, exact primitive top-face parity, 9-sheet determinant audit,
% marked finite-wreath group-level Nielsen/Tietze reduction, and
% canonical nested transfer/idempotent flags.

\subsection{Three-power odd-index towers: actual Schreier relators,
primitive Fox parity, and strict nested Hecke flags}
\label{subsec:dyadic-f135-3power-tower}

The index-three comparison of fixed133--fixed134
is the first stage of a genuinely recursive
odd-index phenomenon.  To test the difference
between a \emph{chosen} finite-index Fox gauge
and a \emph{canonical} higher transfer operation,
we use the tower
\[
K_0=\mathbf Q_2,\qquad
K_r=\mathbf Q_2(\beta_r),\quad
\beta_r^{3^r}=2,\quad
\beta_{r-1}=\beta_r^3,\qquad r\ge1.
\]
Every step has index three,
and the composite degree is
$m_r=3^r$.
On the one hand, the full marked
two-relator source has a recursively
computable, genuinely integral
$m_r$-sheet Schreier--Fox carrier.
Its trivial-coefficient Fox matrix
has one explicit parity obstruction
and no hidden higher integral
elementary divisor.
On the other hand, the canonical
permutation module of the tower
admits a \emph{transversal-independent}
flag of \emph{strictly commuting}
odd-index idempotents.
In the first new nine-sheet stage,
these produce the three integral
ranks $1,2,6$.
The two results give a
concrete extension of the
fixed134 comparison while
distinguishing nested-cover
coherence from the stronger
nonnested or globally preferred
minimal-Fox problem.

Write
\[
\Gamma=G_{\mathbf Q_2},\qquad
H_r=G_{K_r},\qquad
m=3^r,\qquad
\mathscr X_m=H_r\backslash\Gamma .
\]
For this subsection only,
indices of the tame-root sheets
belong to $\mathbf Z/m$.
The polynomial $X^m-2$
is Eisenstein over $\mathbf Q_2$.
The $m$th roots of unity generate
an unramified extension of
degree
\[
f_r=\operatorname{ord}_{m}(2)
=2\cdot3^{r-1},
\]
and the Galois closure
$K_r(\zeta_m)$ has
the tame affine group
\begin{equation}\label{eq:dyadic-f135-tame-affine-group}
\boxed{\quad
D_m=C_m\rtimes C_{f_r},\qquad
\tau:j\longmapsto j+1,\qquad
\sigma:j\longmapsto u_mj,\quad
2u_m\equiv1\pmod m.
\quad}
\end{equation}
Here $\sigma$ has the geometric
Frobenius normalization
$\tau^\sigma=\tau^2$ of
Roe--Turturean
\cite{RoeTurtureanGQ2}.
The permutation action
of $\sigma$ has order
$f_r$ and its pro-$2$
part has order exactly
two:
\begin{equation}\label{eq:dyadic-f135-tame-involution}
\boxed{\qquad
u_m^{3^{r-1}}\equiv-1\pmod m.
\qquad}
\end{equation}

\begin{lemma}[The actual $m$-sheet tame root quotient]
\label{lem:dyadic-f135-tame-affine-root-quotient}
For every $r\ge1$,
$[K_r:\mathbf Q_2]=3^r$,
the Galois closure described
above has group $D_m$,
and the marked wild generators
$x_0,x_1$ map to the identity
in this tame quotient.
In the ordered nine-sheet
case $m=9$, the marked
permutations are explicitly
\begin{equation}\label{eq:dyadic-f135-nine-tame-permutations}
\boxed{\quad
\tau=(0\,1\,2\,3\,4\,5\,6\,7\,8),\qquad
\sigma:j\mapsto5j\bmod9,\qquad
f_2=6.
\quad}
\end{equation}
The intermediate quotient
$\mathscr X_9\to\mathscr X_3$
is $j\mapsto j\bmod3$.
\end{lemma}

\begin{proof}
Eisenstein's criterion
gives the odd degree
$3^r$ and total tame
ramification.
Since $m$ is odd,
the extension generated
by $\zeta_m$ is
unramified, and its
degree is
$\operatorname{ord}_m(2)$.
The fact that $2$ is a
primitive root modulo
each power of three
gives
$\operatorname{ord}_{3^r}(2)
=2\cdot3^{r-1}$.
The finite tame Galois
closure is generated
by the translations of
the roots
$\beta_r\zeta_m^j$
and Frobenius on
$\zeta_m$, giving
\eqref{eq:dyadic-f135-tame-affine-group}.
Its wild inertia is
trivial, so the
marked wild generators
map to the identity.
The congruence
$2^{3^{r-1}}\equiv-1
\pmod{3^r}$
gives
\eqref{eq:dyadic-f135-tame-involution}.
The nine-sheet
formula and reduction
modulo three follow
immediately.
\end{proof}

Let $V_m=\mathbf Z_2^{\,m}$
be the finite coset
permutation lattice.
Denote the permutation
matrices of
\eqref{eq:dyadic-f135-tame-affine-group}
by $\Sigma_m,\Upsilon_m$.
Put
\begin{equation}\label{eq:dyadic-f135-coset-projector-and-augmentation}
\boxed{
P_m=\frac1m\sum_{j=0}^{m-1}\Upsilon_m^j
=\frac1m\mathbf1_m\mathbf1_m^t,
\qquad
Q_m=I_m-P_m.
}
\end{equation}
Both are integral
idempotents over
$\mathbf Z_2$.
They commute with
the full tame affine
permutation action.

\begin{theorem}[Uniform explicit $3^r$-sheet marked Fox differential]
\label{thm:dyadic-f135-uniform-coinduced-fox}
Let
$\mathcal J_m^\bullet$
be the \emph{literal}
two-relator marked
coefficient complex
of the
Roe--Turturean
source evaluated
in the coset
permutation module
$V_m$.
In cochain
degrees $0,1,2$
it has free terms
\[
V_m\longrightarrow
V_m^4\longrightarrow
V_m^2.
\]
Write generator
cochains as
$(a,b,c,d)$ in
the order
$(\sigma,\tau,x_0,x_1)$,
and relation cochains
as (tame,wild).
The complete
integral differential
matrices are
\begin{equation}\label{eq:dyadic-f135-uniform-fox-differentials}
\boxed{
\begin{aligned}
d_m^0(v)&=
\bigl(
(\Sigma_m-I_m)v,
(\Upsilon_m-I_m)v,
0,0\bigr),\\
d_m^1(a,b,c,d)
&=
\left(
\begin{aligned}
&\Sigma_m^{-1}
(\Upsilon_m-I_m)a\\
&\quad+
(\Sigma_m^{-1}-I_m-\Upsilon_m)b,
\end{aligned}\right.\\[-2pt]
&\hspace{3.5em}
\left.
\begin{aligned}
&3P_m b+(4P_m-2I_m)c\\
&\quad+
(\Sigma_m^{-1}-P_m)d
\end{aligned}
\right).
\end{aligned}}
\end{equation}
All coefficients belong to
$\mathbf Z_2$:
the only denominator
is the odd unit $m$.
The matrices satisfy
$d_m^1d_m^0=0$
\emph{exactly}.
For $m=3$ they specialize
to the coinduced
fixed133 Fox row.
For $m=9$ they give an
explicit $18\times36$
matrix, rather than
a merely existential
finite-cover carrier.
\end{theorem}

\begin{proof}
Evaluate each marked
word on the affine
semidirect group
$V_m\rtimes D_m$,
with multiplication
\[
(A,v)(B,w)
=(AB,v+Aw).
\]
The residual tame
permutations have
orders $m$ and
$f_r=2\cdot3^{r-1}$.
For a lift
$(\Upsilon_m,b+c)$
of the product
$x_0\tau$,
the $2$-primary
idempotent power
removes its
odd permutation
part and leaves
the pure translation
\[
(x_0\tau)^{\omega_2}
=(I_m,P_m(b+c)).
\]
The analogous formula
for $x_1\tau$ is
$(I_m,P_m(b+d))$.
Because
$f_r$ has exactly
one factor of two,
the $2$-primary
part of
$(\Sigma_m,a)$
has permutation
of order two.
Its square,
denoted $g_0$ in
the marked wild
words, is a pure
translation.
Therefore it
centralizes all
wild source
translations.
The auxiliary
marked wild words
reduce in this
affine evaluation
to
\[
d_0=P_m(b+c)-c,\qquad
c_0=h_c=1,\qquad
h_0=2c+4d_0.
\]
Finally
$x_1^\sigma$
evaluates to the
translation
$\Sigma_m^{-1}d$.
Consequently the
wild relation
$h_0u_1^{-1}
x_1^\sigma c_0$
has translation
\[
2c+4(P_m(b+c)-c)
-P_m(b+d)+\Sigma_m^{-1}d,
\]
which simplifies
to the second
component of
\eqref{eq:dyadic-f135-uniform-fox-differentials}.
The tame relation
$\sigma^{-1}\tau
\sigma\tau^{-2}$
gives the first
component by the
crossed-derivation
rule.
The ordinary
Fox--crossed-derivation
identity makes this
exactly the
completed marked
coefficient
differential.
Since each relation
is the identity at
the residual coset
permutation
representation,
the principal
coboundary satisfies
$d_m^1d_m^0=0$.
This may also be
checked by the
matrix relation
$\Sigma_m^{-1}
\Upsilon_m\Sigma_m
=\Upsilon_m^2$
and
$P_m\Upsilon_m=P_m$.
\end{proof}

Define the integral
\emph{relative
top-face trace}
on the
$2m$-dimensional
marked relation
space by
\begin{equation}\label{eq:dyadic-f135-uniform-top-face-parity}
\boxed{\quad
\Lambda_m(t,w)
=3\sum_{i=0}^{m-1}t_i
 +\sum_{i=0}^{m-1}w_i,
\qquad (t,w)\in V_m^2 .
\quad}
\end{equation}
It is important
that $\Lambda_m$
is a functional
on the \emph{actual
two-face marked
cover}, not the
fixed132
coefficient-exchange
Laplacian.

\begin{theorem}[Exact analytic image and the unique dyadic Smith factor]
\label{thm:dyadic-f135-analytic-smith-parity}
For every $m=3^r$,
the complete
image of the
marked Fox differential
\eqref{eq:dyadic-f135-uniform-fox-differentials}
is the \emph{exact}
index-two sublattice
\begin{equation}\label{eq:dyadic-f135-image-is-parity-kernel}
\boxed{\qquad
\operatorname{im}d_m^1
=\bigl\{(t,w)\in
\mathbf Z_2^{2m}:
\Lambda_m(t,w)\in2\mathbf Z_2
\bigr\}.
\qquad}
\end{equation}
More precisely, set
\begin{equation}\label{eq:dyadic-f135-discrete-anti-difference}
\boxed{
\mathsf R_m:=
\frac1m\sum_{j=0}^{m-1}
j\,\Upsilon_m^j,
\qquad
(\Upsilon_m-I_m)\mathsf R_m
=Q_m.
}
\end{equation}
For any target
$(t,w)$ satisfying
the parity condition,
the following \emph{literal
integral} source
is a right inverse
to the face
differential on
that target:
\begin{equation}\label{eq:dyadic-f135-exact-face-solver}
\boxed{
\begin{aligned}
a&=\mathsf R_m
 \Sigma_m Q_m t,\\
b&=-P_m t,\\
c&=\tfrac12P_m(w+3t),\\
d&=\Sigma_m Q_m w.
\end{aligned}
\qquad
d_m^1(a,b,c,d)=(t,w).
}
\end{equation}
The factor
$1/2$ in the
\emph{notation}
for $c$ is not
an inversion
of two in the
coefficient ring:
$P_m(w+3t)$ is
already divisible
by two precisely
under the
displayed parity
hypothesis.

Consequently
the \emph{integral
Smith invariants}
of $d_m^1$ are
\begin{equation}\label{eq:dyadic-f135-all-3power-smith}
\boxed{\quad
\operatorname{SNF}_{\mathbf Z_2}(d_m^1)
=(\underbrace{1,\dots,1}_{2m-1},2)
\quad}
\end{equation}
with the remaining
$4m-2m$ source
columns zero.
Equivalently,
$\operatorname{rank}_{\mathbf F_2}
(d_m^1\bmod2)=2m-1$
and
\[
H^2(\mathcal J_m^\bullet)
\cong\mathbf F_2,
\quad
[q]\longmapsto
\Lambda_m(q)\bmod2.
\]
The degree-zero
map has exactly
$m-1$ integral
unit Smith factors,
so the complete
formal cohomology
ledger is
\begin{equation}\label{eq:dyadic-f135-formal-m-cohomology}
\boxed{\quad
H^0(\mathcal J_m^\bullet)=
\mathbf Z_2,\qquad
H^1(\mathcal J_m^\bullet)=
\mathbf Z_2^{m+1},\qquad
H^2(\mathcal J_m^\bullet)=
\mathbf F_2.
\quad}
\end{equation}
These are
\emph{analytically
derived}, and
do not rely on
postulating an
unseen higher
Fox syzygy.
Their
identification
with the actual
field cohomology
is addressed by
the group-level
Schreier theorem
below.
\end{theorem}

\begin{proof}
Write $\mathbf1^t$
for the sum
of coordinates.
Every tame coset
permutation preserves
this functional,
so
\[
\mathbf1^tP_m
=\mathbf1^t,\quad
\mathbf1^tQ_m=0,\quad
\mathbf1^t\Sigma_m^{-1}
=\mathbf1^t,\quad
\mathbf1^t\Upsilon_m
=\mathbf1^t.
\]
Applying
$\Lambda_m$ to
\eqref{eq:dyadic-f135-uniform-fox-differentials}
then gives
\begin{equation}\label{eq:dyadic-f135-face-trace-boundary}
\boxed{\quad
\Lambda_m d_m^1(a,b,c,d)
=2\sum_i c_i.
\quad}
\end{equation}
Thus the
image is contained
in the asserted
parity sublattice.

The finite
geometric-sum
identity
\[
(\Upsilon_m-I_m)
\sum_{j=0}^{m-1}
j\Upsilon_m^j
=
mI_m-\sum_{j=0}^{m-1}
\Upsilon_m^j
\]
proves
\eqref{eq:dyadic-f135-discrete-anti-difference}.
Because $m$ is odd,
$\mathsf R_m$
is integral,
and
$\Upsilon_m-I_m$
is invertible
on the augmentation
summand $Q_mV_m$.
For any target
with even
$\Lambda_m$,
the coefficient
$\tfrac12P_m(w+3t)$
belongs to
$\mathbf Z_2^m$,
since $m$ is odd
and every
coordinate of
the constant
vector $P_m(w+3t)$
is $\Lambda_m(t,w)/m$.
Substitute the
four source
formulas of
\eqref{eq:dyadic-f135-exact-face-solver}.
In the tame row,
\[
\Sigma_m^{-1}
(\Upsilon_m-I_m)
\mathsf R_m\Sigma_mQ_mt
=Q_mt,
\]
and
$(\Sigma_m^{-1}-I_m-\Upsilon_m)
(-P_mt)=P_mt$.
The tame output
is therefore $t$.
In the wild row,
\[
3P_mb=-3P_mt,\quad
(4P_m-2I_m)c
=P_m(w+3t),\quad
(\Sigma_m^{-1}-P_m)d
=Q_mw.
\]
The wild output
is $w$.
This proves the
reverse inclusion
and the exact
image formula.

The parity
functional is
primitive, so
its kernel
modulo two
has index two.
Therefore the
cokernel of
$d_m^1$ is
exactly $\mathbf F_2$,
and the only
nonunit Smith
factor is $2$;
the mod-two
rank is $2m-1$.
For the
degree-zero map,
the $\tau$
component is
$\Upsilon_m-I_m$,
whose restriction
to $Q_mV_m$
is invertible
by
\eqref{eq:dyadic-f135-discrete-anti-difference}.
Its kernel
is exactly the
constant line.
It follows that
$d_m^0$ has
$m-1$ unit
factors and a
saturated image
in $\ker d_m^1$.
The remaining
degree-one
cohomology is
free of rank
$4m-2m-(m-1)=m+1$,
giving
\eqref{eq:dyadic-f135-formal-m-cohomology}.
\end{proof}

\begin{corollary}[The first new nine-sheet determinant certificate]
\label{cor:dyadic-f135-nine-minor-certificate}
At $m=9$ put
\[
\Sigma=\Sigma_9,\quad
\Upsilon=\Upsilon_9,\quad
J_9=\mathbf1_9\mathbf1_9^t.
\]
Multiply the bottom
nine wild-face
rows of $d_9^1$
by the odd unit
$9$.
The resulting
\emph{integer}
$18\times36$
matrix is the
following explicit
block matrix
(each block is
$9\times9$):
\begin{equation}\label{eq:dyadic-f135-nine-block-fox}
\boxed{
\mathscr D_9=
\begin{pmatrix}
\Sigma^{-1}(\Upsilon-I)&
 \Sigma^{-1}-I-\Upsilon&
 0&0\\
0&3J_9&
4J_9-18I&
9\Sigma^{-1}-J_9
\end{pmatrix}.
}
\end{equation}
Order the
36 source
coordinates as
$(\sigma_0,\ldots,\sigma_8,
\tau_0,\ldots,\tau_8,
x_{0,0},\ldots,x_{0,8},
x_{1,0},\ldots,x_{1,8})$,
and the
18 target rows
as nine tame
followed by
nine wild.
Choose the
following
\emph{one-based}
17 source columns,
listed in this order:
\begin{equation}\label{eq:dyadic-f135-nine-pivot-column-set}
\mathcal C_{17}=
(1,2,3,4,5,6,7,8,\,
18,\,
28,29,30,31,32,33,34,35).
\end{equation}
These columns
are the
eight nonterminal
$\sigma$ edges,
the one
non-tree $\tau$
edge, and
the eight
nonterminal
$x_1$ edges;
\emph{none}
is one of the
eight $\tau$
spanning-tree
edges.
Taking the nine
tame rows
$(1,\ldots,9)$
and eight
wild rows
$(11,\ldots,18)$
gives the
exact odd
unit minor
\begin{equation}\label{eq:dyadic-f135-nine-17minor}
\boxed{
\det\mathscr D_9[
(1,\ldots,9,11,\ldots,18),
\mathcal C_{17}]
=-3^{14}.
}
\end{equation}
Adding the final
source column
$27$ (the
$x_{0,8}$ edge)
and the missing
wild row $10$
gives the
full determinant
\begin{equation}\label{eq:dyadic-f135-nine-18minor}
\boxed{
\det\mathscr D_9[
(1,\ldots,18),
(\mathcal C_{17},27)]
=-2\cdot3^{16}.
}
\end{equation}
Thus the selected
\emph{true}
nine-sheet
Fox differential
has 17 unit
directions and
one exact
two-primary
top direction.
After contracting
the eight
$\tau$ tree edges,
the integral
source has
28 generators
and 18 relation
faces; its
augmented
Fox matrix
has Smith factors
\[
(\underbrace{1,\ldots,1}_{17},2).
\]
\end{corollary}

\begin{proof}
The block formula
is
\eqref{eq:dyadic-f135-uniform-fox-differentials}
with $P_9=J_9/9$,
after the declared
odd-unit row
scaling.
The two
integer determinants
are evaluated by
ordinary fraction-free
elimination using
\eqref{eq:dyadic-f135-nine-tame-permutations};
the exact matrices
and both selected
minors are part
of the
independent Sage
audit for this
version.
For the
intrinsic proof
of the Smith
factors without
any determinant
calculation, use
Theorem~\ref{thm:dyadic-f135-analytic-smith-parity}.
The displayed
odd minor additionally
certifies the
\emph{specific}
17 source pivots
used in the
group-level
Tietze construction
below.
\end{proof}

\begin{theorem}[Uniform pro-$2$ Schreier--Tietze reduction for the $3$-power tower]
\label{thm:dyadic-f135-actual-3power-pro2-schreier}
For every $r\ge1$
with $m=3^r$,
the maximal pro-$2$
local Galois group
$P_r=G_{K_r}(2)$
has an
\emph{actual marked}
Schreier source with
\begin{equation}\label{eq:dyadic-f135-general-schreier-counts}
\boxed{\qquad
3m+1\text{ completed pro-$2$
generators},\qquad
2m\text{ indexed relations}.
\qquad}
\end{equation}
Choose the
tree edges
$\tau:0\to1,\ldots,
\tau:m-2\to m-1$.
The
$3m+1$
Schreier
source generators
are
\begin{equation}\label{eq:dyadic-f135-general-schreier-words}
\boxed{
\begin{gathered}
s_i=
\tau^i\sigma
\tau^{-(2i\bmod m)}
\quad(0\le i<m),\\
t=\tau^m,\qquad
w_{i,j}=
\tau^i x_j\tau^{-i}
\quad(0\le i<m,\ j=0,1).
\end{gathered}}
\end{equation}
The $2m$
completed relation
words are defined
without ambiguity
as the actual
Reidemeister--Schreier
rewrites of
$\tau^ir_t\tau^{-i}$
and
$\tau^ir_{\rm wild}\tau^{-i}$
in the
Roe--Turturean
\emph{marked}
profinite presentation,
followed by the
maximal pro-$2$
quotient.
They remain
finite-jet
computable at
every specified
finite pro-$2$
precision.

Exactly $2m-1$
of these
relations are
\emph{primitive}
independent
pro-$2$ Frattini
directions:
all $m$ indexed
tame relations
and any fixed
$m-1$ of the
$m$ indexed
wild relations
can be selected.
Their
linearized
integral Fox
minor is a
unit.
Consequently
one may replace
$2m-1$
source generator
letters by these
$2m-1$
relation words
through a
completed
free pro-$2$
Nielsen
automorphism,
then kill the
corresponding
generator/relation
pairs.
The resulting
\emph{actual}
minimal
presentation is
\begin{equation}\label{eq:dyadic-f135-mplus2-one-relator}
\boxed{
P_r\cong
\langle
y_1,\ldots,y_{m+2}
\mid \mathscr R_r
\rangle_{\mathrm{pro}-2}.
}
\end{equation}
Its relator
$\mathscr R_r$
is the image of
the unique
surviving indexed
wild relation
under the inverse
Nielsen change,
with the
$2m-1$ cancelled
letters set
to one.

For $m=9$,
the explicit
unit minor
\eqref{eq:dyadic-f135-nine-17minor}
certifies the
concrete choice
\emph{all nine tame
relations, wild
relations indexed
$1,\ldots,8$},
and generator
pivots
\[
s_0,\ldots,s_7,\quad
t,\quad
w_{0,1},\ldots,w_{7,1}.
\]
The eleven
retained
source generators
are then
\begin{equation}\label{eq:dyadic-f135-nine-retained-generators}
\boxed{\qquad
s_8,\quad
w_{0,0},\ldots,w_{8,0},
\quad w_{8,1}.
\qquad}
\end{equation}
This is the first
\emph{new}
nine-sheet
group-level
odd-index Fox
reduction
beyond fixed133.

The minimal
relator's exact
closed finite
word in the
retained
generators is
not asserted
to exist.
Its inverse
limit is however
computable
in every finite
pro-$2$ word
quotient.
It represents
the standard
local Demu\v{s}kin
group of
rank $m+2$,
invariant $q=2$,
and full
cyclotomic
orientation image.
In an
\emph{additional}
(normal-form)
Demu\v{s}kin
basis it has
the familiar
rank-odd
representative
\begin{equation}\label{eq:dyadic-f135-standard-3power-normal-form}
y_1^2y_2^4
[y_2,y_3]
[y_4,y_5]\cdots
[y_{m+1},y_{m+2}],
\end{equation}
whose quadratic
Fox cup matrix
is nondegenerate.
This normal-form
change is not
claimed to
coincide with
the marked
Schreier--Nielsen
change above.
\end{theorem}

\begin{proof}
The finite tame
quotient
$D_m$ has an
$m$-point coset
action.
Apply profinite
Reidemeister--Schreier
rewriting to the
marked
four-generator,
two-relator
Roe--Turturean
presentation using
the coset
representatives
$1,\tau,\ldots,
\tau^{m-1}$.
The standard
spanning-tree count
gives
$1+m(4-1)=3m+1$
free generators.
Every global
marked relation
has one lift
at each coset
vertex, hence
there are
$2m$ indexed
relation words,
whose source
coordinates
are precisely
\eqref{eq:dyadic-f135-general-schreier-words}.
The words
containing
$\omega_2$
are interpreted
by their
continuous
profinite powers
before taking
the pro-$2$
quotient.

The marked
wild pro-$2$
normal-closure
condition is
essential.
To see that the
indexed relations
still describe
the \emph{actual}
maximal pro-$2$
subgroup,
test them on an
arbitrary finite
$2$-group $Q$.
A finite quotient
of the subgroup
to $Q$ induces,
by the chosen
$m$ cosets,
a wreath
representation
of $\Gamma$
into
$Q^m\rtimes D_m$,
with every marked
wild generator
in the normal
$2$-group $Q^m$.
Its wild normal
closure is
therefore a
$2$-group, so
the induced
finite quotient
is admissible
for the
Roe--Turturean
marked source.
Conversely,
each finite
$Q$-valued
solution to
the indexed
Schreier relations
produces this
admissible
wreath quotient.
The universal
finite $2$-group
quotient functors
agree, yielding
the \emph{actual}
maximal pro-$2$
presentation
\eqref{eq:dyadic-f135-general-schreier-counts}.
No naked
unmarked
full-group
two-relator
presentation
has been
substituted
at this step.

At the
pro-$2$
Frattini
level, the
abelianized Fox
matrix of this
cover is
identified,
after the
$m-1$ tree-edge
unit cancellations,
with the
completed
coinduced
Fox matrix
\eqref{eq:dyadic-f135-uniform-fox-differentials}.
Its mod-two
rank is $2m-1$
by
Theorem~\ref{thm:dyadic-f135-analytic-smith-parity}.
More precisely,
the $m$ tame
relation rows
span a
primitive
rank-$m$
direct summand:
the $\sigma$
columns supply
the augmentation
directions because
$\Upsilon_m-I_m$
is invertible
on $Q_mV_m$,
and the one
non-tree $\tau$
column supplies
the constant
direction.
Modulo that
tame relation
summand, the
wild $x_1$
block reduces to
$\Sigma_m^{-1}-P_m$,
whose image is
exactly the
$m-1$ dimensional
augmentation module.
Any $m-1$
indexed wild
relation rows
span the dual
of this
augmentation
module.
Thus all
$m$ tame plus
$m-1$ wild
relations
are Frattini
independent;
a suitable
$(2m-1)\times
(2m-1)$
completed Fox
minor is an
odd unit.
The explicit
nine-sheet
instance is
\eqref{eq:dyadic-f135-nine-17minor}.

Keep the
$m+2$
complementary
generator
letters.
The map on the
free pro-$2$
source sending
the selected
$2m-1$
letters to
the selected
relation words
and fixing
the complementary
letters has
invertible
Frattini linear
part.
The pro-$2$
Burnside basis
theorem makes
it surjective;
the finitely
generated free
pro-$2$ group
is Hopfian,
so it is a
continuous
automorphism.
In the new
basis, the
selected
relations are
generator
letters and
can be
killed by
actual
pro-$2$
Tietze operations,
leaving exactly
$m+2$ generators
and one
relation.
This proves
\eqref{eq:dyadic-f135-mplus2-one-relator}
and the
specific
nine-sheet
choice.

Every
pro-$2$ Nielsen
inverse can be
evaluated
in compatible
finite pro-$2$
quotients;
continuous
$\omega_2$
powers of
the marked
source become
finite powers
at each finite
stage.
This proves
the stated
recursive
effectivity.
Finally
$K_r/\mathbf Q_2$
has odd degree,
hence intersects
the cyclotomic
$2$-power
extension
trivially.
It contains
$\mu_2$
but not
$\mu_4$,
so its
canonical
orientation
has full image
$\mathbf Z_2^\times$.
Local
Demu\v{s}kin
theory
\cite{LabuteClassification,NSW}
gives rank
$[K_r:\mathbf Q_2]+2
=m+2$,
$q=2$, and
the additional
standard normal
form
\eqref{eq:dyadic-f135-standard-3power-normal-form}.
The marked
inverse-limit
relator is
not identified
with that
standard word
without its
further
normal-form
change.
\end{proof}

\begin{theorem}[Uniform integral covering-to-minimal Fox chain equivalence]
\label{thm:dyadic-f135-uniform-integral-fox-tietze}
Let $M$ be any
finite projective
complete
coefficient
module for
$H_r$ whose
action factors
through the
maximal pro-$2$
quotient $P_r$.
Let
$C_{\mathrm{cov},m}^\bullet(H_r,M)$
denote the
\emph{actual}
$m$-sheet
marked
Reidemeister--Schreier
cochain complex
with degrees
$(m,4m,2m)$.
Let
\[
C_{\min,r}^\bullet(P_r,M):
\qquad
M\longrightarrow
M^{m+2}\longrightarrow M
\]
be the
completed Fox
complex of
the chosen
minimal
rank-$(m+2)$
one-relator
presentation
\eqref{eq:dyadic-f135-mplus2-one-relator}.
Then the
chosen
Schreier tree
and the
$2m-1$
pro-$2$ Nielsen
pivots construct
a strict integral
chain isomorphism
\begin{equation}\label{eq:dyadic-f135-general-fox-chain-split}
\boxed{
\begin{aligned}
C_{\mathrm{cov},m}^\bullet(H_r,M)
\cong{}&
C_{\min,r}^\bullet(P_r,M)\\
&\oplus
[M^{m-1}\xrightarrow{\,I\,}
 M^{m-1}]_{[0,1]}\\
&\oplus
[M^{2m-1}\xrightarrow{\,I\,}
 M^{2m-1}]_{[1,2]}.
\end{aligned}}
\end{equation}
The degree-wise
ranks match:
\[
(m,4m,2m)
=
(1,m+2,1)
+(m-1,m-1,0)
+(0,2m-1,2m-1).
\]
Since $P_r$
is a local
Demu\v{s}kin
group,
$C_{\min,r}^\bullet$
computes
$R\Gamma(H_r,M)$
and is compatible
with derived
coefficient
base change.
For the new
nine-sheet case
$m=9$,
the explicit
decomposition is
\begin{equation}\label{eq:dyadic-f135-nine-fox-chain-split}
\boxed{
C_{\mathrm{cov},9}^\bullet
\cong
C_{\min,2}^\bullet
\oplus
[M^8\xrightarrow I M^8]_{[0,1]}
\oplus
[M^{17}\xrightarrow I M^{17}]_{[1,2]},
}
\end{equation}
with minimal
degree-wise
ranks $(1,11,1)$.

The entire
comparison is
\emph{recursively
native-explicit}
in each prescribed
finite pro-$2$
word/coefficient
jet, but depends
on the Schreier
tree, ordered
relation pivots,
and inverse
Nielsen basis
map.  This is
a genuine
group-level
integral
Fox syzygy
comparison
for the
selected
$3$-power
radical tower,
not a
presentation-independent
strict gauge
for all
odd-index
local covers.
\end{theorem}

\begin{proof}
The actual
group presentation
and chosen
pro-$2$
Nielsen
automorphism are
Theorem~\ref{thm:dyadic-f135-actual-3power-pro2-schreier}.
The $m-1$
spanning-tree
edges give
literal unit
vertex--edge
Fox pivots,
so they split
$m-1$
contractible
degree-zero--one
pairs.
The remaining
$3m+1$
generator
letters and
$2m$ indexed
relators form
the true
completed
pro-$2$
Schreier
group source.
The selected
$2m-1$
relation words
are a free
pro-$2$
basis replacement.
Its completed
Fox chain
Jacobian is
invertible,
by the Fox
chain rule
and the
inverse Nielsen
automorphism.
In the new
basis each
selected
relation is
literally a
generator,
with Fox
derivative one
in its own
coordinate.
Successive
completed
Fox row/column
operations
therefore split
$2m-1$
contractible
degree-one--two
pairs.
The surviving
relation is
$\mathscr R_r$,
and its
minimal
Fox complex
is the remaining
summand.
All operations
are integral
over
$\mathbf Z_2[[P_r]]$:
no division by
two occurs.
The Demu\v{s}kin
duality
resolution theorem
makes the
one-relator
complex a
length-two
projective
resolution,
so this
source-level
chain isomorphism
computes
continuous
cohomology
for every
coefficient
module
factoring through
$P_r$.
Finite
Frattini
quotients
and continuity
of completed
Fox derivatives
supply the
declared
finite-jet
effectivity.
\end{proof}

\begin{theorem}[Strict nested odd-index Hecke flags]
\label{thm:dyadic-f135-nested-hecke-flags}
Fix $r\ge1$,
$m=3^r$,
and consider
the chain of
subgroups
\[
\Gamma=H_0\supset
H_1\supset\cdots\supset H_r,
\qquad
[H_{i-1}:H_i]=3.
\]
Let
\[
\pi_i:
\mathscr X_m=H_r\backslash\Gamma
\longrightarrow
H_i\backslash\Gamma,
\qquad
0\le i\le r
\]
be the
canonical
$\Gamma$-equivariant
coset projections.
The fiber
of $\pi_i$
has cardinal
$3^{r-i}$.
For an arbitrary
complete local
$\mathbf Z_2$-algebra
$A$ put
$V=A[\mathscr X_m]$.
Define the
\emph{finite
fiber-average}
operator
\begin{equation}\label{eq:dyadic-f135-flag-fiber-averages}
\boxed{\quad
E_i=
\frac1{3^{r-i}}
\,\pi_i^*\pi_{i,!}:
V\longrightarrow V,
\qquad 0\le i\le r,
\quad}
\end{equation}
where $\pi_{i,!}$
sums the
coefficients
over each
fiber and
$\pi_i^*$
copies a
coefficient
back to
that fiber.
Then all $E_i$
are
\emph{integral
$\Gamma$-equivariant
idempotents}, and
for every
$i,j$,
\begin{equation}\label{eq:dyadic-f135-strict-flag-composition}
\boxed{\qquad
E_iE_j=E_jE_i
=E_{\min(i,j)}.
\qquad}
\end{equation}
The difference
operators
\begin{equation}\label{eq:dyadic-f135-flag-primitive-layers}
\boxed{
D_0=E_0,\qquad
D_i=E_i-E_{i-1}
\ (1\le i\le r)
}
\end{equation}
are pairwise
orthogonal
idempotents,
sum to one,
and have
\emph{exact}
projective
ranks
\begin{equation}\label{eq:dyadic-f135-flag-layer-ranks}
\boxed{\qquad
\operatorname{rank}_A D_0=1,
\qquad
\operatorname{rank}_A D_i
=3^i-3^{i-1}
\quad(i\ge1).
\qquad}
\end{equation}
In the
new nine-sheet
case $r=2$,
with coset
coordinate
$j\in\mathbf Z/9$,
\begin{equation}\label{eq:dyadic-f135-nine-explicit-flag}
\boxed{
\begin{aligned}
E_0&=\tfrac19
\mathbf1_9\mathbf1_9^t,\\
(E_1)_{jk}
&=\begin{cases}
1/3,&j\equiv k\pmod3,\\
0,&j\not\equiv k\pmod3,
\end{cases}\\
E_2&=I_9,\qquad
(D_0,D_1,D_2)
=(E_0,E_1-E_0,I_9-E_1),\\
(\operatorname{rank}D_0,
\operatorname{rank}D_1,
\operatorname{rank}D_2)
&=(1,2,6).
\end{aligned}}
\end{equation}
All identities
are literal
equalities of
$9\times9$
integral matrices;
their matrix
entries contain
no factor $1/2$.

Let $B$ be any
finite projective
complete $A$-module
with continuous
$\Gamma$-action.
There is an
intrinsic
coinduction
identification
\begin{equation}\label{eq:dyadic-f135-coinduced-complete-flag}
\operatorname{Coind}_{H_r}^\Gamma
(\operatorname{Res}_{H_r}^\Gamma B)
\simeq B\otimes_A V.
\end{equation}
Thus
$1_B\otimes D_i$
defines an exact
$\Gamma$-equivariant
coefficient
decomposition.
The operators
are \emph{strict
cochain maps}
on both
the continuous
bar complex
and the finite
marked
two-relator
Fox complex:
they commute
literally with
every differential.
Whenever the
restriction
$B|_{H_r}$
has pro-$2$
image, the
group-level
completion of
Theorem~\ref{thm:dyadic-f135-uniform-integral-fox-tietze}
and Shapiro give
the \emph{sectorwise}
derived
decomposition
\begin{equation}\label{eq:dyadic-f135-derived-layer-decomposition}
\boxed{\quad
R\Gamma(H_r,B|_{H_r})
\simeq
\bigoplus_{i=0}^r
R\Gamma\bigl(
\Gamma,
B\otimes_A \operatorname{im}D_i
\bigr).
\quad}
\end{equation}
The $D_0$
summand is
$R\Gamma(\Gamma,B)$,
and the
composite
odd-index
restriction--corestriction
projector is
exactly the
constant-sheet
$D_0$.

These
identities provide
\emph{strict
associativity}
of the nested
odd-index
Hecke projector
system:
reparenthesizing
a tower of
fiber sums
changes no
map and
requires no
higher homotopy
filler.
They are
independent of
the chosen
Schreier
transversals;
a change of
transversal
conjugates all
coordinate
matrices by the
same
coefficient-transport
change of basis.
They do
\emph{not}
claim the same
strict
composition for
crossing,
nonnested
Sylow
preimages.
\end{theorem}

\begin{proof}
Because the
$\pi_i$ are
maps of finite
$\Gamma$-sets,
the fiber sums
$\pi_{i,!}$
and pullbacks
$\pi_i^*$
are
$\Gamma$-equivariant.
Every fiber
has odd
cardinality,
so its inverse
exists in $A$.
The composite
$\pi_{i,!}\pi_i^*$
is multiplication
by $3^{r-i}$;
therefore
$E_i$ is an
idempotent.
For $i\le j$,
each fiber of
$\pi_i$ is a
disjoint union
of equal-size
fibers of
$\pi_j$.
Averaging
inside the
smaller
fibers and
then in the
larger one
is exactly
the larger
average,
regardless of
the order.
This proves
\eqref{eq:dyadic-f135-strict-flag-composition}.
The $D_i$
therefore form
pairwise
orthogonal
idempotents
whose direct
sum is identity.
Their images
are the new
layers added
by the
successive
fiber-constant
functions;
their ranks
are the
differences
$3^i-3^{i-1}$.
The nine-sheet
formulas follow
by writing
$\pi_1(j)=j\bmod3$.

For a
$\Gamma$-module
$B$, the
standard
finite-index
coinduction
trivialization
uses the
already present
$\Gamma$-action
on $B$ to
identify
the fiber
module with
$B\otimes A[\mathscr X_m]$.
Under this
identification
the $\Gamma$
action is
diagonal.
The coefficient
projectors
$1_B\otimes D_i$
commute with
this action.
Consequently
they induce
strict
cochain
projections
on the
continuous
bar complex.
Completed
Fox cochain
matrices are
finite sums
of coefficient
group-action
operators,
including
limits of
profinite powers;
they therefore
commute with
these same
projectors.
This gives
the strict
finite
Fox splitting.

Shapiro's
lemma
identifies
the complete
coinduced
bar complex
with
$R\Gamma(H_r,B|_{H_r})$.
If $B|_{H_r}$
has pro-$2$
image,
Theorem~\ref{thm:dyadic-f135-uniform-integral-fox-tietze}
identifies the
literal $m$-sheet
marked complex
with an
actual finite
Fox resolution
of $P_r$,
so the finite
direct summands
also compute
the intrinsic
derived
cohomology.
The constant
summand equals
$B$, hence its
cohomology is
$R\Gamma(\Gamma,B)$,
and the
integral average
$E_0$ is
restriction
followed by
normalized
corestriction
under Shapiro.
The strict
composition
identity proves
the asserted
nested higher
coherence.
A change of
coset sections
only changes
the chosen
coordinate
trivialization,
not the
intrinsic
fiber maps.
\end{proof}

\begin{corollary}[Full nine-sheet integral arithmetic cohomology ledger]
\label{cor:dyadic-f135-nine-arithmetic-ledger}
For the concrete
two-step tower
\[
K_0=\mathbf Q_2,\quad
K_1=\mathbf Q_2(2^{1/3}),\quad
K_2=\mathbf Q_2(2^{1/9}),
\]
the complete
low-degree
integral
arithmetic
cohomology
for the trivial
and cyclotomic
rank-one
lattices is
\begin{equation}\label{eq:dyadic-f135-nine-h0h1h2-ledgers}
\boxed{
\begin{array}{c|ccc|ccc}
&\multicolumn{3}{c|}{\mathbf Z_2}
&\multicolumn{3}{c}{\mathbf Z_2(1)}\\
& H^0&H^1&H^2&H^0&H^1&H^2\\ \hline
K_0&
\mathbf Z_2&
\mathbf Z_2^2&
\mathbf F_2&
0&\mathbf Z_2^2\oplus\mathbf F_2&\mathbf Z_2\\
K_1&
\mathbf Z_2&
\mathbf Z_2^4&
\mathbf F_2&
0&\mathbf Z_2^4\oplus\mathbf F_2&\mathbf Z_2\\
K_2&
\mathbf Z_2&
\mathbf Z_2^{10}&
\mathbf F_2&
0&\mathbf Z_2^{10}\oplus\mathbf F_2&\mathbf Z_2
\end{array}}
\end{equation}
The
strict
rank-$(1,2,6)$
projector flag
selects
exactly
\[
\begin{aligned}
&\operatorname{rank}_{\mathbf Z_2}
H^1(K_0,\cdot)=2,\\
&\operatorname{rank}_{\mathbf Z_2}
\text{new layer }K_1/K_0=2,\\
&\operatorname{rank}_{\mathbf Z_2}
\text{new layer }K_2/K_1=6,
\end{aligned}
\]
for both twists.
Every $H^0$
or $H^2$
class and
every
$\mathbf F_2$
torsion class
of
$H^1(-,\mathbf Z_2(1))$
lies in the
\emph{oldest}
constant-sheet
summand.
The new
rank-two and
rank-six
summands have
cohomology
concentrated
in degree one:
\begin{equation}\label{eq:dyadic-f135-new-layer-degree-one-only}
\boxed{
R\Gamma(\Gamma,\operatorname{im}D_1(\nu))
\simeq\mathbf Z_2^2[-1],
\quad
R\Gamma(\Gamma,\operatorname{im}D_2(\nu))
\simeq\mathbf Z_2^6[-1],
\quad
\nu\in\{0,1\},
}
\end{equation}
where $(\nu)$
denotes tensoring
with the
$\nu$th
cyclotomic twist.
In particular
the rank-six
newest layer
is the actual
integral
Shapiro-induced
fiber augmentation:
\begin{equation}\label{eq:dyadic-f135-new-six-layer-induction}
\boxed{\quad
\operatorname{im}D_2
\simeq
\operatorname{Ind}_{H_1}^\Gamma
\ker\bigl(
\mathbf Z_2[H_1/H_2]
\xrightarrow{\rm sum}
\mathbf Z_2
\bigr).
\quad}
\end{equation}
\end{corollary}

\begin{proof}
Every $K_i$
is a finite
dyadic field
of odd degree
$3^i$,
hence contains
exactly $\mu_2$
among its
$2$-power
roots of unity.
Local class
field theory,
Kummer theory
and integral
Tate duality
give
\[
H^1(G_{K_i},\mathbf Z_2)
\simeq\mathbf Z_2^{3^i+1},
\quad
H^2(G_{K_i},\mathbf Z_2)
\simeq\mathbf F_2,
\]
and
\[
H^1(G_{K_i},\mathbf Z_2(1))
\simeq\mathbf Z_2^{3^i+1}
\oplus\mathbf F_2,
\quad
H^2(G_{K_i},\mathbf Z_2(1))
\simeq\mathbf Z_2 .
\]
These give
\eqref{eq:dyadic-f135-nine-h0h1h2-ledgers}.
Odd-index
restriction is
split injective
in every
cohomological
degree, with
left inverse
normalized
corestriction.
Since the
displayed
$H^0$ and
$H^2$ groups
have constant
rank/order,
restriction
is an
isomorphism
on them;
for $\mathbf Z_2(1)$,
restriction
on $H^2$
is multiplication
by $3$ or $9$
under normalized
local Tate
invariants,
hence still
a $2$-adic
unit.
Likewise
the unique
$\mathbf F_2$
torsion in
cyclotomic
$H^1$
is preserved
by odd-index
restriction.
Thus the two
new projective
layers have
no degree-zero,
degree-two
or degree-one
torsion
cohomology.
Their degree-one
free ranks
are the
differences
$4-2=2$ and
$10-4=6$.
Since the
bar/Fox
carriers are
perfect over
the discrete
valuation ring
$\mathbf Z_2$,
these ledgers
give the
derived
identifications
\eqref{eq:dyadic-f135-new-layer-degree-one-only}.
Finally
$\operatorname{im}D_2$
consists of
functions with
zero sum
inside each
fiber of
$\mathscr X_9\to\mathscr X_3$.
This is exactly
the induction
from $H_1$
of the
two-dimensional
fiber augmentation
module of
$H_1/H_2$,
proving
\eqref{eq:dyadic-f135-new-six-layer-induction}.
\end{proof}

\begin{warning}[Why nested odd-index coherence does not settle crossing Sylows]
\label{warn:dyadic-f135-crossing-sylow-firewall}
The composition
law
\eqref{eq:dyadic-f135-strict-flag-composition}
uses \emph{nested}
odd-index coset
maps, so the
fiber sizes
at every stage
are units in
$\mathbf Z_2$.
For two
\emph{distinct}
Sylow-$2$
subgroups
$S,S'$ inside
the same
$S_3$ quotient,
their index-three
preimages
$H,H'$ are
not nested;
their
intersection is
the inverse
image of
$S\cap S'=\{1\}$,
hence
\[
[\Gamma:H\cap H']=6.
\]
A universal
constant-fiber
average across
this common
refinement
would require
division by
two and is
\emph{not}
an integral
$\mathbf Z_2$
operator.
Already for
$C_2$ the
augmentation sequence
\[
0\to
\ker(\varepsilon:
\mathbf Z_2[C_2]\to
\mathbf Z_2)
\to
\mathbf Z_2[C_2]
\xrightarrow{\varepsilon}
\mathbf Z_2\to0
\]
does not admit a
$C_2$-equivariant
splitting:
the invariant
constant vector
has augmentation
two, not one.
Thus strict
odd-index
nested projector
flags cannot
be extended
\emph{by the same
averaging formula}
to arbitrary
nonnested
Sylow intersections.

Nor does the
uniform
$3$-power
Schreier--Nielsen
algorithm
select one
canonical
marked minimal
Fox basis
independently
of every choice
of transversal,
local field
or residual
quotient.
It supplies
genuine
group-level
finite-jet
chain
comparisons
throughout
the chosen
$3$-power
tame radical
tower and
strict
coherence of
its nested
transfer
idempotents.
The full
arithmetic
Cartan
derived
faithful atlas
remains
unconditional;
global
presentation-independent
strict
gauge coherence
across
\emph{crossing}
residual
coverings
remains a
stronger
separate
problem.
\end{warning}
